% Generated draft: requires mathematical and source review.
\documentclass[11pt]{article}
\usepackage[margin=25mm]{geometry}
\usepackage{fontspec}
\setmainfont{Latin Modern Roman}
\usepackage{amsmath,amssymb,mathtools,unicode-math}
\setmathfont{Latin Modern Math}
\usepackage{enumitem,needspace,xurl,hyperref}
\setlength{\parindent}{0pt}
\setlength{\parskip}{5pt}
\setlength{\emergencystretch}{3em}
\newcommand{\N}{\mathbb{N}}
\newcommand{\Z}{\mathbb{Z}}
\newcommand{\Q}{\mathbb{Q}}
\newcommand{\R}{\mathbb{R}}
\newcommand{\C}{\mathbb{C}}
\newcommand{\F}{\mathbb{F}}
\newcommand{\A}{\mathbb{A}}
\newcommand{\PP}{\mathbb{P}}
\newcommand{\E}{\mathbb{E}}
\newcommand{\eps}{\varepsilon}
\newcommand{\dd}{\,\mathrm{d}}
\begin{document}
\section*{UnsolvedMath ID: 30003006}
\textbf{Source problem code:} catalogue entry
\subsection{Definitions and mathematical statement}
The supplied record collects several questions, but does not define the operations needed to formulate all of them precisely. For standard notation only, in a finite-dimensional Euclidean space \(\R^n\), a convex body means a compact convex set with nonempty interior. For \(0\in A\), define its polar set by
\[
A^\circ=\{y\in\R^n:\langle x,y\rangle\leq 1\text{ for every }x\in A\}.
\]
Write \(D=\{x\in\R^n:\|x\|_2\leq 1\}\), define Minkowski addition by \(A+B=\{a+b:a\in A,\ b\in B\}\), and let \(|K|\) denote \(n\)-dimensional Lebesgue measure when applicable. The source does not specify the dimension or its range, whether symmetry is required, or whether containing the origin means \(0\in A\) or \(0\in\operatorname{int}A\). These are unresolved hypotheses: if the origin lies on the boundary, the polar set need not be a convex body.

Subject to these unresolved domain qualifications, the record asks the following questions about convex bodies \(A,B\) containing the origin:
\begin{enumerate}
\item Does
\[
G(A,B)=D\quad\Longrightarrow\quad B=A^\circ
\]
hold for every admissible pair \((A,B)\)?
\item Does
\[
A+B=A^\circ+B^\circ\quad\Longrightarrow\quad B=A^\circ
\]
hold for every admissible pair \((A,B)\)?
\item Must a convex body \(P\) be an ellipsoid if it is representable as \(G(A,B)\) for every admissible \(A\)? A natural reading of the representation condition is
\[
\text{for every admissible }A\text{ there exists an admissible }B
\text{ such that }P=G(A,B),
\]
with \(B\) allowed to depend on \(A\). This quantifier interpretation requires confirmation; the record does not explicitly quantify \(B\) or specify the admissible class of \(P\).
\item Does \(G_p(A,B)=G_1(A,B)\) hold? The definition of \(G_p\), the allowed values of \(p\), and the intended quantification over \(p,A,B\) are absent.
\item Does the upper elliptic envelope satisfy \(E_u(A,B)=G(A,B)\)? The construction of \(E_u\) is absent.
\item Under which conditions on \(A,B\), and possibly the ambient dimension, does
\[
|G(A,B)|^2\geq |A|\,|B|
\]
hold? The record asks for identifiable conditions and does not assert this inequality without additional hypotheses.
\item Does an exponential map exist, and does the referenced homogeneity identity hold? The record supplies neither the domain, codomain, nor required properties of the exponential map, and does not state the homogeneity identity.
\end{enumerate}
The symbols \(G\) and \(G_p\) are described only as geometric means of convex bodies. No construction or characterizing axioms are supplied, so these questions cannot yet be interpreted as fully specified mathematical propositions. In particular, no choice of geometric mean, elliptic envelope, exponential map, or homogeneity identity is made here.
\subsection{Short English statement}
The record asks whether geometric means of convex bodies characterize polar pairs and ellipsoids, agree with parameterized means and elliptic envelopes, satisfy a volume inequality, and admit an exponential map and a homogeneity identity. Essential definitions and some hypotheses are missing, so this is a statement-only draft requiring human mathematical review. The supplied open label is a dated curation assessment, not an independently verified status.
\subsection{Sources}
\begin{itemize}
\item Local registry: \texttt{data/unsolvedmath/problems.json}, with \texttt{catalog/problems.json} designated as the supplementary registry; permanent UnsolvedMath ID \(30003006\), exact source problem code \texttt{catalogue entry}, title \emph{Structural Properties of Geometric Means of Convex Bodies}. These identifiers are taken from the supplied complete record; no registry file was independently inspected.
\item Supplied report citation: \emph{Convex Geometry and its Applications} (2016), Oberwolfach Reports, pp.\ 3179--3269, DOI \texttt{10.4171/owr/2015/56}. Supplied source URL and literature reference, labeled ``Convex Geometry and its Applications report'': \url{https://doi.org/10.4171/owr/2015/56}. No individual proposer is supplied. The record gives proposed year 2015, while the citation labels the report 2016; both are retained.
\item Supplied literature assessment, checked on exactly \texttt{2026-08-21}: status ``open.'' The supplied assessment says that the report lists general questions about polar characterization, geometric means, elliptic envelopes, and \(p\)-means; that later work gives special-case results; and that no general theorem resolves all extracted clauses. This is supplied dated curation, not independently verified. No references for those later special-case results are supplied. The background explicitly requests independent verification of time-sensitive status and cited claims.
\item The supplied statement status is \texttt{corrected\_verified}; its accompanying qualification says that only safe notation, typography, and TeX-delimiter cleanup was applied, with hypotheses, quantifiers, and subquestions preserved. This label does not establish mathematical correctness or resolve the missing definitions.
\end{itemize}
\end{document}
